A further reduction allows ordinary locally nilpotent ideals in the
upper idempotent rings. For a unital ring $A$, let $\operatorname{Lev}(A)$
be its largest locally nilpotent two-sided ideal: every finite subset
of this ideal generates a nilpotent, possibly nonunital, subring.

\begin{theorem}[a compatible locally nilpotent quotient]
\label{thm:upper-levitzki-quotient}
Let $T=(Q,F)$ be any unital inverse semiring. Define $H_q$ by the tests
\[
 x\in H_q\quad\Longleftrightarrow\quad
 t_{pqr,e}(uxv)\in\operatorname{Lev}(F(e))
\]
for every $u\in F(p)$, $v\in F(r)$ and $e\in U(Q)$ with $pqr\leq e$.
Then $H$ is the largest coefficient ideal whose fiber at each
$e\in U(Q)$ is locally nilpotent. It contains $K^{\mathrm{id}}$, and
$T/H=(Q,F/H)$ has no nonzero coefficient ideal with that same property.

For every finite idempotent support matrix $E$, the kernel of
$A_E(T)\to A_E(T/H)$ is locally nilpotent. Extension of scalars gives
a direct-sum-preserving bijection of finitely generated projective
left-semimodule isomorphism classes, and hence
\[
 V(T)\simeq V(T/H),\qquad K_0(T)\simeq K_0(T/H).
\]
For every $m\geq1$, this coefficient ideal commutes with finite matrices:
\[
 H(M_m(T))_P=\prod_{i,j}H(T)_{P_{ij}}
 \qquad(P\in M_m(Q)).
\]
Consequently $M_m(T)/H(M_m(T))\simeq M_m(T/H(T))$ canonically.
\end{theorem}

The nilpotence assertion is local inside each ordinary matrix fiber;
there is no bound depending only on matrix size. For an ordinary ring
$T=R$, the quotient is the familiar $R/\operatorname{Lev}(R)$, so the
ordinary radical and nil-ideal lifting are known inputs. The new
coefficient condition concerns compatible ideals. It does not make
every individual upper idempotent ring radical-free. Exact supports
are retained, and no equivalence of full projective categories,
arbitrary-map functoriality or arbitrary Morita-equivalence assertion
is made. The unrestricted support and amplitude classifications
remain to be evaluated after this further coefficient reduction.


Proposition~\ref{prop:upper-transition-radical} below evaluates this
same ideal using only ordinary-ring radicals and upward transitions.
Every coefficient group uses two-by-two upper-idempotent fiber rings;
for finite $Q$ the intersection is finite, even when its coefficient
groups are infinite. This is a formula, not a general radical algorithm.
